\section{Trabalhos Relacionados}
\label{sec:trabalhos_relacionados}

\subsection{Modelos compartimentais e doenças transmitidas por mosquitos}

O modelo SIR de \citeonline{kermack1927} divide a população em suscetíveis, infectados e recuperados e
mostra que um surto só acontece quando o número básico de reprodução, $R_0$, passa de 1. Para doenças
transmitidas por mosquitos, a base teórica é a de Ross e Macdonald, que representa a transmissão como um
ciclo entre humanos e mosquitos e mostra que ela é proporcional à densidade de mosquitos por pessoa e muito
sensível à sobrevivência do mosquito adulto \cite{smith2012}. Essa teoria orienta até hoje o controle
vetorial.

\subsection{Modelos hospedeiro e vetor para a chikungunya}

\citeonline{yakob2013} ajustaram um modelo determinístico com humanos e mosquitos à grande epidemia da ilha
da Reunião, em 2006, e estimaram $R_0$ de 4,1. Mesmo simples, o modelo reproduziu o pico e o tamanho final
da epidemia, e a análise de sensibilidade mostrou que a latência humana e o tempo de incubação no mosquito
pesam muito no ajuste, o que justifica representá-los separadamente. \citeonline{manore2014} compararam
dengue e chikungunya em \textit{Aedes aegypti} e \textit{Aedes albopictus} com modelos SEIR-SEI,
reunindo os parâmetros da literatura para os dois vetores. O trabalho também deixa clara uma distinção
que afeta a leitura de qualquer estimativa de $R_0$: o valor calculado pela matriz de próxima geração é a
raiz quadrada do número de reprodução de humano para humano.

Os parâmetros biológicos têm fontes próprias. A revisão de \citeonline{rudolph2014} estima a incubação da
chikungunya em humanos em cerca de 3 dias. \citeonline{christofferson2014} mostram que a incubação no
mosquito tem média próxima de 8 dias para a linhagem ECSA e que usar o tempo mais curto em que o vírus
aparece na saliva superestima a velocidade da epidemia. A revisão de \citeonline{haider2020} reúne
estimativas de $R_0$ de surtos transmitidos por \textit{Aedes}, com média de 3,4, mas não traz nenhuma
estimativa para o Brasil.

\subsection{A chikungunya no Brasil e no Recife}

A disseminação do vírus pelas Américas a partir de 2013 é descrita por \citeonline{weaver2015}. A
linhagem ECSA entrou no Brasil por Feira de Santana, na Bahia, em 2014 \cite{faria2016}, e o país teve
sete ondas epidêmicas até 2022, com recorrência em vários estados \cite{desouza2023}. Na Região
Metropolitana do Recife, um estudo com pacientes febris de uma unidade de pronto atendimento em Paulista
detectou o vírus em junho de 2015 e mostrou a chikungunya substituindo a zika a partir de agosto daquele
ano \cite{magalhaes2017}.

Os inquéritos sorológicos mostram que a fração infectada varia muito entre lugares. Depois das primeiras
epidemias, foram encontrados 57\% de soropositivos em setores de Feira de Santana \cite{dias2018}, 11,8\%
numa comunidade de Salvador \cite{anjos2020} e 35,7\% no Recife, com 22,3\% no estrato socioeconômico alto e
cerca de 39\% nos estratos intermediário e baixo \cite{braga2023}. Uma meta-análise estima 29\% em média
nas Américas \cite{skalinski2023}. Essa variação é um argumento direto contra a hipótese de mistura
homogênea.

A hipótese de imunidade permanente, ao contrário, tem bom apoio. Anticorpos neutralizantes persistem por
décadas \cite{nitatpattana2014}, protegem contra novas infecções \cite{srikiatkhachorn2016} e reagem com
genótipos diferentes, o que confirma que o vírus tem um só sorotipo \cite{auerswald2018}. É uma diferença
importante em relação à dengue, que tem quatro tipos de vírus \cite{cdc2026}.

\subsection{Monitoramento e controle}

O InfoDengue estima a transmissão de arboviroses em centenas de municípios brasileiros a partir das
notificações, corrigindo o atraso de notificação e calculando o número de reprodução efetivo
\cite{codeco2018}. É uma ferramenta estatística, que não representa o mosquito nem a imunidade acumulada.
Sobre o controle, a revisão sistemática de \citeonline{bowman2016} encontrou poucos ensaios randomizados
que mediram o efeito do controle vetorial sobre a incidência de dengue e concluiu que a evidência é fraca
e pouco padronizada, embora o manejo ambiental comunitário com tampas nos reservatórios tenha mostrado
efeito. Por isso a eficácia de uma campanha deve ser tratada como incerta e explorada em cenários.

\subsection{O que esta proposta acrescenta}

Os modelos hospedeiro e vetor da literatura foram ajustados a outros lugares, e as ferramentas de
monitoramento usadas no Brasil não são mecanísticas. Esta proposta junta as duas coisas para o Recife: um
modelo SEIR-SEI calibrado com a série local de notificações e com o inquérito sorológico da própria cidade,
capaz de estimar a imunidade remanescente e de comparar estratégias de controle.
